%% ma: L11-B14-0008
%% lop: 11
%% chuong: 4
%% bai: 14
%% dang: 11.14.2
%% loai: TN4
%% muc: NB
%% dap_an: "A"
%% nguon: "tai-lieu-hoc-tap-mon-toan-lop-11-theo-chuong-trinh-giao-duc-pho-thong-2018.pdf (Lớp Toán Thầy Duy, Phạm Lê Duy, THPT Chu Văn An), Chương 4 Bài 5, Phần 1, Câu 3 (tr.523)"
%% kien_thuc: "Lớp 11 Bài 14 (phép chiếu song song)"
%% trang_thai: chinh-thuc
%% ghi_chu: "Đã duyệt 03/11/2026."
\begin{ex}
Qua phép chiếu song song, tính chất nào \textbf{không} được bảo toàn?
\choice{\True Chéo nhau.}{Đồng quy.}{Thẳng hàng.}{Song song.}
\loigiai{\sloppy Phép chiếu song song bảo toàn tính thẳng hàng, tính song song và tính đồng quy (hai đường cắt nhau có hình chiếu cắt nhau hoặc trùng nhau), nhưng hai đường chéo nhau có thể có hình chiếu song song hoặc cắt nhau, nên tính chéo nhau không được bảo toàn.}
\end{ex}
